problem:  
Let $\{(M_i^n,g_i)\}$ be a sequence of closed Riemannian $n$‑manifolds with $\mathrm{sec}_{g_i}\ge -1$ and $\mathrm{diam}(M_i)\le 1$, collapsing in the Gromov–Hausdorff sense to a compact Alexandrov space $(X^k,d_X)$ with $k<n$.  

For each $x\in X$, let $T_xX$ denote the tangent cone, a metric cone over the space of directions $\Sigma_x$ (an Alexandrov space of curvature $\ge1$).  Define two points $x,y\in X$ to be **cone‑equivalent** if there is an isometry $T_xX\cong T_yX$ of metric cones.  Let $\mathcal{C}(X)$ be the set of cone‑equivalence classes of points in $X$.  

**Question:**  
For fixed $n$ and uniform lower sectional curvature bound (say $\mathrm{sec}\ge -1$), is there a number $N(n)<\infty$ such that for every collapse $(M_i^n,g_i)\to X^k$ as above, the corresponding limit $X$ satisfies $|\mathcal{C}(X)|\le N(n)$?  
Equivalently, does there exist a uniform *finiteness of tangent‑cone types* for Alexandrov spaces arising as limits of $n$‑dimensional manifolds with $\mathrm{sec}\ge -1$ and bounded diameter?

Why is it a "good" problem:  
This problem avoids undefined objects while retaining deep geometric substance.  It probes whether collapsing with a lower sectional curvature bound produces only finitely many distinct local geometries (tangent cone types) in the limit.  Such a finiteness theorem would represent a major extension of known results: in two‑sided curvature bounds, the local models are nilpotent Lie group quotients, so a finite list of local types exists; under one‑sided bounds, no such result is known.  

A positive answer would impose powerful structural constraints on Alexandrov limits; a negative one would reveal new sources of singular behavior.  The question is simple, precise, and yet fundamentally open—it sits at the interface of metric geometry, topology, and geometric analysis, and resolving it would advance the understanding of how lower curvature bounds control analytic and topological complexity in the collapsing regime.  Its clear statement, potential universality, and profound implications make it a “good” research‑level problem.